\section{协方差和相关系数}

前面介绍的数学期望和方差，无论是否为二维变量，最终考察的总是一维变量的数字特征。
本节介绍描述两个变量之间关系的数字特征——协方差和相关系数。

本节要点：
\begin{itemize}
    \item 掌握协方差的概念；
    \item 掌握相关系数的概念；
    \item 深刻理解相关系数的由来及几何意义。
\end{itemize}

%============================================================
\subsection{协方差的概念}

\begin{definition}[协方差]
设二维随机变量$\left( X,Y \right) $，若存在数学期望
\[
E\left[ \left( X-EX \right) \left( Y-EY \right) \right]
\]
则称之为{\bf $X$和$Y$的协方差}，记作$\mathrm{Cov}\left( X,Y \right) $，即：
\begin{align*}
\mathrm{Cov}\left( X,Y \right) :&=\sum_i{\sum_j{\left( x_i-EX \right) \cdot \left( y_j-EY \right) \cdot p_{ij}}} \\
\mathrm{Cov}\left( X,Y \right) :&=\iint\limits_{\mathbb{R} \times \mathbb{R}}{\left( x-EX \right) \cdot \left( y-EY \right) \cdot f\left( x,y \right) \cdot d\sigma}, \quad d\sigma =dxdy
\end{align*}
\end{definition}

由数学期望和方差的性质可得便捷计算公式：
\begin{align*}
\mathrm{Cov}\left( X,Y \right) &=E\left[ \left( X-EX \right) \left( Y-EY \right) \right] \\
&=E\left[ XY-X\cdot EY-Y\cdot EX+EX\cdot EY \right] \\
&=E\left( XY \right) -EY\cdot E\left( X \right) -EX\cdot E\left( Y \right) +E\left( EX\cdot EY \right) \\
&=E\left( XY \right) -EX\cdot EY
\end{align*}
特别地，当$X$和$Y$相互独立时$E\left( XY \right) =E\left( X \right) \cdot E\left( Y \right) $，有：
\[
\mathrm{Cov}\left( X,Y \right) =E\left( X \right) \cdot E\left( Y \right) -EX\cdot EY=0
\]

\begin{tcolorbox}
$\mathrm{Cov}$是“Covariance”的缩写。同样定义式更多的是为了理解概念，实际中通常使用便捷公式。
\end{tcolorbox}

%============================================================
\subsection{协方差的性质}

{\bf 性质1}：$\mathrm{Cov}\left( X,X \right) =DX$

{\bf 性质2}：$\mathrm{Cov}\left( X,c \right) =0$

{\bf 性质3}：设$\left( X,Y \right) $为二维随机变量，

$D\left( X\pm Y \right) =DX+DY\pm 2\mathrm{Cov}\left( X,Y \right) $

{\bf 性质4}：设$\left( X,Y \right) $为二维随机变量，$\mathrm{Cov}\left( X,Y \right) =\mathrm{Cov}\left( Y,X \right) $

{\bf 性质5}：设$\left( X,Y \right) $为二维随机变量，$\left[ \mathrm{Cov}\left( X,Y \right) \right] ^2\leqslant DX\cdot DY$

{\bf 性质6}：设$\left( X,Y \right) $为二维随机变量，$\mathrm{Cov}\left( aX,bY \right) =ab\mathrm{Cov}\left( X,Y \right) $

{\bf 性质7}：设$\left( X_1,X_2,Y \right) $为三维随机变量，

$\mathrm{Cov}\left( X_1\pm X_2,Y \right) =\mathrm{Cov}\left( X_1,Y \right) \pm \mathrm{Cov}\left( X_2,Y \right) $

~

性质1，用各自的便捷公式即可证明：
\[
\mathrm{Cov}\left( X,X \right) =E\left( X^2 \right) -\left( EX \right) ^2=DX
\]

性质3，用各自的便捷公式也可证明：
\begin{align*}
D\left( X\pm Y \right) &=\left[ E\left( X^2 \right) \pm 2E\left( XY \right) +E\left( Y^2 \right) \right] - \\
&\left[ \left( EX \right) ^2\pm 2\left( EX \right) \left( EY \right) +\left( EY \right) ^2 \right] \\
&=DX+DY\pm 2\left[ E\left( XY \right) +\left( EX \right) \left( EY \right) \right] \\
&=DX+DY\pm 2\mathrm{Cov}\left( X,Y \right)
\end{align*}

\begin{tcolorbox}
上述两个性质联系了方差和协方差，特别是性质3，可配合方差性质：
\[
\text{若}X,Y\text{相互独立}\Rightarrow D\left( X\pm Y \right) =D\left( X \right) +D\left( Y \right)
\]
进行理解。
\end{tcolorbox}

性质5，是理解相关系数的关键，证明如下：
\begin{align*}
&\because \mathrm{Cov}\left( X,Y \right) =E\left[ \left( X-EX \right) \left( Y-EY \right) \right] \\
&\because \left[ E\left( XY \right) \right] ^2\leqslant E\left( X^2 \right) \cdot E\left( Y^2 \right) \\
&\begin{aligned}
	\therefore \left[ \mathrm{Cov}\left( X,Y \right) \right] ^2&=\left[ E\left[ \left( X-EX \right) \left( Y-EY \right) \right] \right] ^2\\
	&\leqslant E\left[ \left( X-EX \right) ^2 \right] \cdot E\left[ \left( Y-EY \right) ^2 \right] =DX\cdot DY\\
\end{aligned}
\end{align*}

%============================================================
\subsection{协方差的统计意义}
从协方差的定义入手：
\[
\mathrm{Cov}\left( X,Y \right) :=\sum_i{\sum_j{\left( x_i-EX \right) \cdot \left( y_j-EY \right) \cdot p_{ij}}}
\]

\begin{itemize}
    \item $\left( x_i-EX \right) \cdot \left( y_j-EY \right) $：表示样本点对均值构成的矩形面积，且1、3象限取正，2、4象限取负；
    \item $\sum_i{\sum_j{\cdot p_{ij}}}$：表示对这些面积加权和。
\end{itemize}
可见：
\begin{itemize}
    \item 当$\mathrm{Cov}\left( X,Y \right) >0$时，表示1、3象限居多，统计上，对于一个样本，$X$取大值时，可能$Y$也是大值；
    \item 反之，当$\mathrm{Cov}\left( X,Y \right) <0$时，统计上，对于一个样本，$X$取大值时，$Y$反倒可能是小值。
\end{itemize}
协方差就表示了$X$和$Y$在分布上有的这个关系，但协方差没有去除$X,Y$各自的方差，下面的相关系数除去了方差的影响，更能精准地刻画$X$和$Y$在分布上的关系。

%============================================================
\subsection{相关系数和不相关的概念}

\begin{definition}[相关系数]
设二维随机变量$\left( X,Y \right) $，若$DX>0,DY>0$，则称$\frac{\mathrm{Cov}\left( X,Y \right)}{\sqrt{DX}\cdot \sqrt{DY}}$为{\bf 二维随机变量$\left( X,Y \right) $的相关系数}，记作$\rho _{XY}$或$\rho $，即：
\[
\rho _{XY}:=\frac{\mathrm{Cov}\left( X,Y \right)}{\sqrt{DX}\cdot \sqrt{DY}}
\]
\end{definition}

由各自的便捷公式可得相关系数的便捷公式：
\[
\rho _{XY}=\frac{\mathrm{Cov}\left( X,Y \right)}{\sqrt{DX}\cdot \sqrt{DY}}=\frac{E\left( XY \right) -EX\cdot EY}{\sqrt{E\left( X^2 \right) -\left( EX \right) ^2}\cdot \sqrt{E\left( Y^2 \right) -\left( EY \right) ^2}}
\]

\begin{definition}[不相关]
对于二维变量$\left( X,Y \right) $，若$\rho _{XY}=0$，则称它们{\bf 不相关}。
\end{definition}

不相关指的是$X,Y$没有线性关系，不排除有其他关系，如平方关系，三角函数关系等，所以不相关是一个比独立要弱的要求。
$X,Y$不相关并不一定独立，而独立必然不相关。

\begin{theorem}
若$X,Y$独立且$\rho _{XY}$存在，则$X,Y$不相关。
\end{theorem}

简单来说即，独立$\Rightarrow $不相关，还是要注意反命题不成立。

\begin{theorem}
若二维变量$\left( X,Y \right) $存在$\rho _{XY}$，则下列命题等价：
\begin{itemize}
    \item $X,Y$不相关；
    \item $\rho _{XY}=0$
    \item $\mathrm{Cov}\left( X,Y \right) =0$
    \item $E\left( XY \right) =EX\cdot EY$
    \item $D\left( X\pm Y \right) =D\left( X \right) +D\left( Y \right) $
\end{itemize}
\end{theorem}

%============================================================
\subsection{相关系数的性质}

{\bf 性质1}：$\rho _{XY}\in \left[ -1,1 \right] $

{\bf 性质2}：$\rho _{XY}=\pm 1$的充要条件是$X$和$Y$呈线性关系，即$\exists a\ne 0,b$使得$P\left\{ Y=aX+b \right\} =1$，且此时有$a>0\Leftrightarrow \rho _{XY}=1,a<0\Leftrightarrow \rho _{XY}=-1$。

{\bf 性质3}：$\forall a\ne 0,\forall b\ne 0$，有
\[
\rho _{\left( aX \right) \left( bY \right)}=\begin{cases}
	\rho _{XY},&		ab>0\\
	-\rho _{XY},&		ab<0\\
\end{cases}
\]

性质1，由相关系数的定义和协方差的性质7可证。

性质2在下面说明。

%============================================================
\subsection{相关系数的由来}

以二维离散变量$\left( X,Y \right) $为例说明，首先我们定义一个函数：
\[
Z=z\left( X,Y \right) =\left[ Y-\left( aX+b \right) \right] ^2
\]
我们求它的数学期望：
\[
EZ=E\left[ z\left( X,Y \right) \right] =\sum_i{\sum_j{\left[ Y-\left( aX+b \right) \right] ^2\cdot p_{ij}}}
\]
若$\left( X,Y \right) $的分布如下左图，我们用圆圈表示$p_{ij}$，圆圈越大表示$p_{ij}$越大，一共7个 ，呈一直线排列。
可见，只有当$X,Y$取值满足$y=ax+b$时，有$p_{ij}\ne 0$，于是$EZ=0$。
若$\left( X,Y \right) $的分布如下中图，4个不在一条直线上的点使得$EZ\ne 0$。
同样，如下右图，$EZ\ne 0$，而且$EZ$会变得更大。

\begin{figure}[h]
\centering
\includegraphics[height=2.2cm]{4.11.png}
\end{figure}

相关系数考察的就是这个关系。

首先引入线性代数的理论，令 $\hat{Y}=aX+b$，即$Y$在$X$上的投影：
\begin{align*}
e\left( a,b \right) &=E\left\{ \left[ Y-\left( aX+b \right) \right] ^2 \right\} \\
&=E\left( Y^2 \right) +a^2E\left( X^2 \right) +b^2-2aE\left( XY \right) -2bE\left( Y \right) +2abE\left( X \right)
\end{align*}

然后我们考察$X$需要做出怎么样的调整使得和$Y$的这个差异最小，即考察$e\left( a,b \right) $的极值。首先用一阶偏导求极值点$\left( a_0,b_0 \right) $：
\begin{align*}
&\begin{cases}
	\frac{\partial e}{\partial a}=2aE\left( X^2 \right) -2E\left( XY \right) +2bE\left( X \right) =0\\
	\frac{\partial e}{\partial b}=2b-2E\left( Y \right) +2aE\left( X \right) =0\\
\end{cases} \\
&\Rightarrow \begin{cases}
	a_0=\frac{\mathrm{Cov}\left( X,Y \right)}{D\left( X \right)}\\
	b_0=E\left( Y \right) -\frac{\mathrm{Cov}\left( X,Y \right)}{D\left( X \right)}\cdot E\left( X \right)\\
\end{cases}
\end{align*}
但这只是极值点，极大值还是极小值还是鞍点，需要考虑$e\left( a,b \right) $的Hessian矩阵的特征值：
\begin{align*}
&H\left( e \right) =\left[ \begin{matrix}
	\frac{\partial ^2e}{\partial a^2}&		\frac{\partial ^2e}{\partial a\partial b}\\
	\frac{\partial ^2e}{\partial b\partial a}&		\frac{\partial ^2e}{\partial b^2}\\
\end{matrix} \right] =\left[ \begin{matrix}
	2E\left( X^2 \right)&		2E\left( X \right)\\
	2E\left( X \right)&		2\\
\end{matrix} \right] \\
&\det \left[ H\left( e \right) \right] =4E\left( X^2 \right) -4\left[ 2E\left( X \right) \right] ^2=4D\left( X \right) >0
\end{align*}
再加上$\frac{\partial ^2e}{\partial a^2}=2E\left( X^2 \right) >0$，所以$\left( a_0,b_0 \right) $是极小值点，$e\left( a,b \right) $在该点有最小值：
\[
\left. e\left( a,b \right) \right|_{\left( a_0,b_0 \right)}=\left[ 1-\left( \frac{\mathrm{Cov}\left( X,Y \right)}{\sqrt{D\left( X \right)}\cdot \sqrt{D\left( Y \right)}} \right) ^2 \right] D\left( Y \right)
\]

上式衡量了$X$和$Y$在分布上的差异量，即 “努力地做出调整”后，$Y$和$X$的差异量。
若要考察一个已知的$Y$对$X_1$和$X_2$的线性关系谁大谁小，则需要排除$Y$的因素，于是我们可取前半部分因子
\[
\left[ 1-\left( \frac{\mathrm{Cov}\left( X_?,Y \right)}{\sqrt{D\left( X_? \right)}\cdot \sqrt{D\left( Y \right)}} \right) ^2 \right]
\]
该因子越小，说明差异量越小，和$X_?$的契合度越高。反之取平方项部分，记作$\rho _{XY}$
\[
\rho _{X_?Y}=\frac{\mathrm{Cov}\left( X_?,Y \right)}{\sqrt{D\left( X_? \right)}\cdot \sqrt{D\left( Y \right)}}
\]
$\rho _{X_?Y}$越大，说明和$X_?$的契合度越高。
也就是说，对于一个给定的$Y$和一堆$X_1,X_2,\cdots X_n$，我们可以用$\rho _{XY}$判断$Y$和某个$X_i$在分布上是最契合的，此时，$X_i$只需要做出调整$a_0X_i+b_0$就可以将差异减小到最小的程度。

对于现实的指导是，如果$\rho _{XY}$大到一定程度，我们就可以用关于$X$的一个简单的函数
\[
g\left( X \right) =\frac{\mathrm{Cov}\left( X,Y \right)}{D\left( X \right)}X+E\left( Y \right) -\frac{\mathrm{Cov}\left( X,Y \right)}{D\left( X \right)}\cdot E\left( X \right)
\]
代替$Y$考察分布。
比如对于正态分布，有很多很好很简单的性质，我们也有很多统计手段和公式。
倘若 “很像”正态分布，那我们就可以用正态分布代替$Y$，进而可以通过公式用标准分布计算，整个问题就变得简单明了。

%============================================================
\subsection{相关的几何意义}

$\rho _{XY}$表示的相关性如下图。
可见，当$\rho _{XY}=\pm 1$时，$Y\ne aX+b$的地方概率为0，$f\left( x,y \right) $将由一张三维曲面退化成一个二维平面。

\begin{figure}[h]
\centering
\includegraphics[height=4.5cm]{4.12.png}
\end{figure}

以二维正态分布为例：
\begin{align*}
&\left( X,Y \right) \sim N\left( \mu _1,\mu _2,{\sigma _1}^2,{\sigma _2}^2,\rho \right) \\
&f\left( x,y \right) =\frac{1}{2\pi \sigma _1\sigma _2\sqrt{1-\rho ^2}}e^{-\frac{1}{2\left( 1-\rho ^2 \right)}\left[ \frac{\left( x-\mu _1 \right) ^2}{{\sigma _1}^2}-2\rho \frac{\left( x-\mu _1 \right) \left( y-\mu _2 \right)}{\sigma _1\sigma _2}+\frac{\left( y-\mu _2 \right) ^2}{{\sigma _2}^2} \right]} \\
&\rho \in \left( -1,1 \right)
\end{align*}
$\rho $表示了$X,Y$的相关性，当$\rho =0.99$时，$f\left( x,y \right) $几乎成为一个平面。

\begin{figure}[h]
\centering
\includegraphics[height=8cm]{4.13.png}
\end{figure}

\begin{tcolorbox}
可见，相关描述了密度函数的扁平性，相关系数就是这个扁平性的量化指标，$\rho _{XY}$越大，曲面$f\left( x,y \right) $越趋于平面。
\end{tcolorbox}

%============================================================
\subsection{协方差和相关系数的线性代数意义}

参考线性代数中向量的柯西—施瓦茨不等式的表述。若$\boldsymbol{x},\boldsymbol{y}$为内积空间$V$中的两个向量，则必有：
\[
\left| \left< \boldsymbol{x},\boldsymbol{y} \right> \right|\leqslant \sqrt{\left< \boldsymbol{x},\boldsymbol{x} \right>}\sqrt{\left< \boldsymbol{y},\boldsymbol{y} \right>}
\]
当且仅当$\boldsymbol{x}=0$或$\boldsymbol{y}=0$或$\boldsymbol{x}=\boldsymbol{y}=0$或$\boldsymbol{x},\boldsymbol{y}$线性相关时等号成立。
于是当$\boldsymbol{x}\ne \boldsymbol{y}$时，我们可以用柯西—施瓦茨不等式定义它们的夹角：
\[
\cos \theta =\frac{\left< \boldsymbol{x},\boldsymbol{y} \right>}{\sqrt{\left< \boldsymbol{x},\boldsymbol{x} \right>}\sqrt{\left< \boldsymbol{y},\boldsymbol{y} \right>}}
\]
对比观察得到，协方差就是两个向量的内积，相关系数就是它们的夹角的余弦：
\[
\rho _{XY}=\frac{\mathrm{Cov}\left( X,Y \right)}{\sqrt{DX}\cdot \sqrt{DY}}=\frac{\mathrm{Cov}\left( X,Y \right)}{\sqrt{\mathrm{Cov}\left( X,X \right)}\cdot \sqrt{\mathrm{Cov}\left( Y,Y \right)}}
\]
各部分对应关系如下：
\[
\begin{array}{l}
	\mathrm{Cov}\left( , \right)\\
	X\\
	Y\\
	DX\\
	DY\\
	\mathrm{Cov}\left( X,Y \right)\\
	\rho _{XY}\\
\end{array} \quad \leftrightarrow \quad \begin{array}{l}
	\left< , \right>\\
	\boldsymbol{x}\\
	\boldsymbol{y}\\
	\left< \boldsymbol{x},\boldsymbol{x} \right>\\
	\left< \boldsymbol{y},\boldsymbol{y} \right>\\
	\left< \boldsymbol{x},\boldsymbol{y} \right>\\
	\cos \theta\\
\end{array}
\]

更进一步，将一维随机变量的分布视为向量，则所有分布的集合可以视作向量空间$V$，任取两种分布$X,Y$用向量记号记作$\boldsymbol{x},\boldsymbol{y}$，我们采用协方差作为操作，即$\left< , \right> =\mathrm{Cov}\left( , \right) $，必然满足内积的4个要求：
\begin{enumerate}[label=\Roman*.]
    \item 正定性，$\left< \boldsymbol{x},\boldsymbol{x} \right> \geqslant 0$，当且仅当$\boldsymbol{x} =0$（即常数分布）时等号成立；
    \item $\left< \alpha \boldsymbol{x},\boldsymbol{y} \right> =\alpha \left< \boldsymbol{x},\boldsymbol{y} \right> $；
    \item 交换律，$\left< \boldsymbol{x},\boldsymbol{y} \right> =\left< \boldsymbol{y},\boldsymbol{x} \right> $；
    \item 分配律，$\left< \boldsymbol{x}+\boldsymbol{y},\boldsymbol{z} \right> =\left< \boldsymbol{x},\boldsymbol{z} \right> +\left< \boldsymbol{y},\boldsymbol{z} \right> $。
\end{enumerate}
其中I、II、III这3条从数学期望的定义式就可得
IV证明如下：
\begin{align*}
&\mathrm{Cov}\left( \left( X+Y \right) ,Z \right) =E\left[ \left( X+Y-E\left( X+Y \right) \right) \left( Z-EZ \right) \right] \\
&=E\left[ \left( X+Y+EX+EY \right) \left( Z-EZ \right) \right] \\
&=E\left( XZ-XEZ+YZ-YEZ+ZEX-EXEZ+ZWY-EZEY \right) \\
&=E\left( XZ \right) +E\left( YZ \right) -E\left( EXEZ \right) -E\left( EZEY \right) \\
&=\mathrm{Cov}\left( X,Z \right) +\mathrm{Cov}\left( Y,Z \right)
\end{align*}

\begin{tcolorbox}
可见，一维随机变量连同协方差运算构成了内积空间，常数分布是零向量，$\rho _{XY}$表示了两个向量的夹角的余弦。
\end{tcolorbox}

%============================================================
\subsection{方差与范数}

更“线性代数化地说”，方差是变量的一个范数的算术平方$DX\leftrightarrow \left\| \boldsymbol{x} \right\| ^2$。范数的要求如下：
\begin{itemize}
    \item 正定性，$\left\| \boldsymbol{x} \right\| \geqslant 0$，当且仅当$\boldsymbol{x}=0$时，等号成立；
    \item $\left\| \alpha \boldsymbol{x} \right\| =\left| \alpha \right|\left\| \boldsymbol{x} \right\| $；
    \item 三角不等式，$\left\| \boldsymbol{x}+\boldsymbol{y} \right\| \leqslant \left\| \boldsymbol{x} \right\| +\left\| \boldsymbol{y} \right\| $。
\end{itemize}
对应方差的性质：
\begin{itemize}
    \item $D\left( X \right) \geqslant 0$，当且仅当$P\left\{ X=EX \right\} =1$时等号成立。
    \item $D\left( kX \right) =k^2\cdot D\left( X \right) $
    \item $D\left( X+Y \right) =DX+DY+2\mathrm{Cov}\left( X,Y \right) $
\end{itemize}

%============================================================
\subsection{例}

\begin{example}
继续第三章中的例子，若5个产品中有2个次品，每次抽取一个，取两次，记$X$表示第一次取到次品的个数，$Y$表示第二次取到次品的个数，讨论：
\begin{enumerate}
    \item 不放回抽取的情况下，
    \item 放回抽取的情况下，
\end{enumerate}
这两种情况下，$X,Y$的协方差。
\end{example}

解：

由题干可得，$X,Y$取值只可能是0或1。

1. 不放回抽取的情况下的分布律$p_{ij}$：
\[
\begin{matrix}
	&		0&		1&		p_{\cdot j}\\
	0&		\frac{3}{10}&		\frac{3}{10}&		\frac{3}{5}\\
	1&		\frac{3}{10}&		\frac{1}{10}&		\frac{2}{5}\\
	p_{i\cdot}&		\frac{3}{5}&		\frac{2}{5}&		1\\
\end{matrix}
\]
\begin{align*}
&\because \begin{cases}
	E\left( XY \right) =0\cdot 0\cdot \frac{3}{10}+0\cdot 1\cdot \frac{3}{10}+1\cdot 0\cdot \frac{3}{10}+1\cdot 1\cdot \frac{1}{10}=\frac{1}{10}\\
	E\left( X \right) =0\cdot \frac{3}{5}+1\cdot \frac{2}{5}=\frac{2}{5}\\
	E\left( Y \right) =0\cdot \frac{3}{5}+1\cdot \frac{2}{5}=\frac{2}{5}\\
\end{cases} \\
&\therefore \mathrm{Cov}\left( X,Y \right) =E\left( XY \right) -E\left( X \right) \cdot E\left( Y \right) =\frac{1}{10}-\frac{4}{25}=-\frac{3}{50}
\end{align*}

2. 放回抽取的情况下的分布律$p_{ij}$：
\[
\begin{matrix}
	&		0&		1&		p_{\cdot j}\\
	0&		\frac{9}{25}&		\frac{6}{25}&		\frac{3}{5}\\
	1&		\frac{6}{25}&		\frac{4}{25}&		\frac{2}{5}\\
	p_{i\cdot}&		\frac{3}{5}&		\frac{2}{5}&		1\\
\end{matrix}
\]
\begin{align*}
&\because \begin{cases}
	E\left( XY \right) =0\cdot 0\cdot \frac{9}{25}+0\cdot 1\cdot \frac{6}{25}+1\cdot 0\cdot \frac{6}{25}+1\cdot 1\cdot \frac{4}{25}=\frac{4}{25}\\
	E\left( X \right) =0\cdot \frac{3}{5}+1\cdot \frac{2}{5}=\frac{2}{5}\\
	E\left( Y \right) =0\cdot \frac{3}{5}+1\cdot \frac{2}{5}=\frac{2}{5}\\
\end{cases} \\
&\therefore \mathrm{Cov}\left( X,Y \right) =E\left( XY \right) -E\left( X \right) \cdot E\left( Y \right) =\frac{4}{25}-\frac{4}{25}=0
\end{align*}

\begin{tcolorbox}
从该题解读协方差的意义。在(2)放回的情况下，从事实角度出发，$X,Y$是独立的，既然独立就没有任何关系，包括线性关系，自然$\mathrm{Cov}\left( X,Y \right) =0$。
在(1)不放回的情况下，第一次抽中次品时，是不利于再次抽中次品这个事件的发生的。
\end{tcolorbox}




